\documentclass[10pt,a4paper]{amsart}
\usepackage{amsmath,amssymb,amsthm,graphicx}
\usepackage{verbatim}
\newenvironment{reference}{\comment}{\endcomment}
%\newenvironment{reference}{}{}
\newenvironment{slogan}{\comment}{\endcomment}
\newenvironment{history}{\comment}{\endcomment}

% For commutative diagrams we use Xy-pic
\usepackage[all]{xy}

% We use 2cell for 2-commutative diagrams.
\xyoption{2cell}
% Theorem environments.
%
\theoremstyle{plain}
\newtheorem{theorem}[subsection]{Theorem}
\newtheorem{proposition}[subsection]{Proposition}
\newtheorem{lemma}[subsection]{Lemma}

\theoremstyle{definition}
\newtheorem{definition}[subsection]{Definition}
\newtheorem{example}[subsection]{Example}
\newtheorem{exercise}[subsection]{Exercise}
\newtheorem{situation}[subsection]{Situation}

\theoremstyle{remark}
\newtheorem{remark}[subsection]{Remark}
\newtheorem{remarks}[subsection]{Remarks}

\numberwithin{equation}{subsection}

% Macros
%
\def\lim{\mathop{\mathrm{lim}}\nolimits}
\def\colim{\mathop{\mathrm{colim}}\nolimits}
\def\Spec{\mathop{\mathrm{Spec}}}
\def\Hom{\mathop{\mathrm{Hom}}\nolimits}
\def\Ext{\mathop{\mathrm{Ext}}\nolimits}
\def\SheafHom{\mathop{\mathcal{H}\!\mathit{om}}\nolimits}
\def\SheafExt{\mathop{\mathcal{E}\!\mathit{xt}}\nolimits}
\def\Sch{\mathit{Sch}}
\def\Mor{\mathop{\mathrm{Mor}}\nolimits}
\def\Ob{\mathop{\mathrm{Ob}}\nolimits}
\def\Sh{\mathop{\mathit{Sh}}\nolimits}
\def\NL{\mathop{N\!L}\nolimits}
\def\CH{\mathop{\mathrm{CH}}\nolimits}
\def\proetale{{pro\text{-}\acute{e}tale}}
\def\etale{{\acute{e}tale}}
\def\QCoh{\mathit{QCoh}}
\def\Ker{\mathop{\mathrm{Ker}}}
\def\Im{\mathop{\mathrm{Im}}}
\def\Coker{\mathop{\mathrm{Coker}}}
\def\Coim{\mathop{\mathrm{Coim}}}

% Boxtimes
%
\DeclareMathSymbol{\boxtimes}{\mathbin}{AMSa}{"02}

%
% Macros for moduli stacks/spaces
%
\def\QCohstack{\mathcal{QC}\!\mathit{oh}}
\def\Cohstack{\mathcal{C}\!\mathit{oh}}
\def\Spacesstack{\mathcal{S}\!\mathit{paces}}
\def\Quotfunctor{\mathrm{Quot}}
\def\Hilbfunctor{\mathrm{Hilb}}
\def\Curvesstack{\mathcal{C}\!\mathit{urves}}
\def\Polarizedstack{\mathcal{P}\!\mathit{olarized}}
\def\Complexesstack{\mathcal{C}\!\mathit{omplexes}}
% \Pic is the operator that assigns to X its picard group, usage \Pic(X)
% \Picardstack_{X/B} denotes the Picard stack of X over B
% \Picardfunctor_{X/B} denotes the Picard functor of X over B
\def\Pic{\mathop{\mathrm{Pic}}\nolimits}
\def\Picardstack{\mathcal{P}\!\mathit{ic}}
\def\Picardfunctor{\mathrm{Pic}}
\def\Deformationcategory{\mathcal{D}\!\mathit{ef}}
\usepackage{cleveref}
\crefname{originallemma}{original lemma}{original lemmas}
\Crefname{originallemma}{Original lemma}{Original lemmas}
\makeatletter
\newlabel{lemma-tensor-product-pullback}{{03EL}{source}{}{sourcepointer-1}{}}
\newlabel{lemma-tensor-product-pullback@cref}{{[originallemma][0][]03EL}{[1][0][]source}{}{}{}}
\makeatother
\crefname{originalreference}{original reference}{original references}
\Crefname{originalreference}{Original reference}{Original references}
\makeatletter
\newlabel{sites-lemma-point-morphism-sites}{{05V1}{source}{}{sourcepointer-2}{}}
\newlabel{sites-lemma-point-morphism-sites@cref}{{[originalreference][0][]05V1}{[1][0][]source}{}{}{}}
\makeatother
\crefname{originalreference}{original reference}{original references}
\Crefname{originalreference}{Original reference}{Original references}
\makeatletter
\newlabel{sites-lemma-point-morphism-topoi}{{05V2}{source}{}{sourcepointer-3}{}}
\newlabel{sites-lemma-point-morphism-topoi@cref}{{[originalreference][0][]05V2}{[1][0][]source}{}{}{}}
\makeatother
\title{Stalks of pullback modules at paired points: (f\textasciicircum{}*\textbackslash{}mathcal\{F\})\_p = \textbackslash{}mathcal\{F\}\_q \textbackslash{}otimes\_\{\textbackslash{}mathcal\{O\}\_\{\textbackslash{}mathcal\{D\},q\}\} \textbackslash{}mathcal\{O\}\_\{\textbackslash{}mathcal\{C\},p\}}
\author{The Stacks Project Authors}
\date{stacks-project:sites-modules.tex:local-snapshot-20261012; GNU FDL 1.2 or later}
\begin{document}
\maketitle
\noindent\textit{Editorial note.} All delivered mathematical text is the author's own, copied verbatim from the pinned Stacks Project chapter file sites-modules.tex; no statement or proof step was rewritten, reordered or removed.\par\medskip
\noindent\textit{Editorial note.} The counted span cites one same-chapter lemma and two Sites lemmas, recorded in assembly.external\_references with canonical tags as navigation pointers.\par\medskip
\noindent\textit{Allowed standard prerequisites (editorial).}\begin{itemize}
\item Stalks commute with tensor products. Stacks Project, Modules on Sites, Lemma lemma-tensor-product-pullback. Supplies that applying p\textasciicircum{}\{-1\} to a tensor product of modules gives the tensor product of the stalks.
\item Stalks of pullbacks under a morphism of sites/topoi. Stacks Project, Sites, Lemmas sites-lemma-point-morphism-sites and sites-lemma-point-morphism-topoi. Supplies the identification of the stalk of f\textasciicircum{}\{-1\}\textbackslash{}mathcal\{F\} at p with the stalk of \textbackslash{}mathcal\{F\} at q.
\end{itemize}\medskip
\ifdefined\hypertarget\hypertarget{sourcepointer-1}{}\fi
\noindent\textit{Source reference.} Original lemma 03EL: sources/stacks-chow/tags.txt, source lines 4460--4460. Cross-chapter standard prerequisite cited by the author inside the counted span; it is used as a pointer and its own proof is not part of this item.\par\medskip
\ifdefined\hypertarget\hypertarget{sourcepointer-2}{}\fi
\noindent\textit{Source reference.} Original reference 05V1: sources/stacks-chow/tags.txt, source lines 7535--7535. Cross-chapter standard prerequisite cited by the author inside the counted span; it is used as a pointer and its own proof is not part of this item.\par\medskip
\ifdefined\hypertarget\hypertarget{sourcepointer-3}{}\fi
\noindent\textit{Source reference.} Original reference 05V2: sources/stacks-chow/tags.txt, source lines 7536--7536. Cross-chapter standard prerequisite cited by the author inside the counted span; it is used as a pointer and its own proof is not part of this item.\par\medskip
\begin{lemma}
\label{lemma-pullback-stalk}
Let
$(f, f^\sharp) :
(\Sh(\mathcal{C}), \mathcal{O}_\mathcal{C})
\to
(\Sh(\mathcal{D}), \mathcal{O}_\mathcal{D})$
be a morphism of ringed topoi or ringed sites.
Let $p$ be a point of $\mathcal{C}$ or $\Sh(\mathcal{C})$
and set $q = f \circ p$. Then
$$
(f^*\mathcal{F})_p =
\mathcal{F}_q \otimes_{\mathcal{O}_{\mathcal{D}, q}}
\mathcal{O}_{\mathcal{C}, p}
$$
for any $\mathcal{O}_\mathcal{D}$-module $\mathcal{F}$.
\end{lemma}

\begin{proof}
We have
$$
f^*\mathcal{F} =
f^{-1}\mathcal{F} \otimes_{f^{-1}\mathcal{O}_\mathcal{D}}
\mathcal{O}_\mathcal{C}
$$
by definition. Since taking stalks at $p$ (i.e., applying
$p^{-1}$) commutes with $\otimes$ by
Lemma \ref{lemma-tensor-product-pullback}
we win by the relation between the stalk of pullbacks at $p$
and stalks at $q$ explained in
Sites, Lemma \ref{sites-lemma-point-morphism-sites} or
Sites, Lemma \ref{sites-lemma-point-morphism-topoi}.
\end{proof}

\end{document}
